\documentclass[11pt]{article}
\usepackage[utf8]{inputenc}
\usepackage{amsmath,amssymb}
\usepackage{hyperref}
\title{Scalable Green Chemistry Sustainable Synthesis for Large-Scale Settings}
\author{Anonymous}
\date{}
\begin{document}
\maketitle

\begin{abstract}
This paper studies Green Chemistry Sustainable Synthesis. Grounded in a real, citable literature \cite{ref1,ref2,ref3,ref4,ref5,ref6,ref7,ref8},
we propose a focused method and report \textbf{illustrative} experiments.
All references below are real and verifiable (OpenAlex/arXiv); the reported
numbers are simulated to illustrate intended behaviour.
\end{abstract}

\section{Introduction}
Green Chemistry Sustainable Synthesis is an active research area. This work targets a concrete gap identified from
the recent literature and contributes a method together with an evaluation protocol.

\section{Related Work (real literature)}
The following works, retrieved from public bibliographic APIs, frame our study:
\begin{itemize}
\item Gallezot (2011) — "Conversion of biomass to selected chemical products", Chemical Society Reviews (cited 2530).
\item Narayanan et al. (2010) — "Biological synthesis of metal nanoparticles by microbes", Advances in Colloid and Interface Science (cited 1897).
\item Graham et al. (2003) — "Legumes: Importance and Constraints to Greater Use", PLANT PHYSIOLOGY (cited 1783).
\item Sheldon et al. (2017) — "Role of Biocatalysis in Sustainable Chemistry", Chemical Reviews (cited 1737).
\item Sheldon (2005) — "Green solvents for sustainable organic synthesis: state of the art", Green Chemistry (cited 1694).
\item Mohanty (2005) — "Natural Fibers, Biopolymers, and Biocomposites" (cited 1603).
\end{itemize}

\section{Method}
We decompose the problem across shards with a communication-efficient aggregation rule that keeps overhead sublinear in scale. We instantiate this idea for Green Chemistry Sustainable Synthesis and analyse its cost/quality trade-off.

\section{Experiments (Illustrative)}
\begin{center}
\begin{tabular}{lcc}
System & Primary Metric & Efficiency \\ \hline
Strong baseline & 0.812 & 1.00$\times$ \\
Ablation & 0.828 & 0.92$\times$ \\
\textbf{Ours} & \textbf{0.851} & \textbf{0.71$\times$}
\end{tabular}
\end{center}
\emph{Metrics simulated to illustrate the intended effect; not measured.}

\section{Conclusion}
We connected a real literature to a concrete method for Green Chemistry Sustainable Synthesis. Future work will
validate the illustrative results with real experiments.

\begin{thebibliography}{9}
\bibitem{ref1} P. Gallezot. Conversion of biomass to selected chemical products. \emph{Chemical Society Reviews}, 2011. DOI:10.1039/c1cs15147a
\bibitem{ref2} Kannan Badri Narayanan, N. Sakthivel. Biological synthesis of metal nanoparticles by microbes. \emph{Advances in Colloid and Interface Science}, 2010. DOI:10.1016/j.cis.2010.02.001
\bibitem{ref3} Peter Graham, Carroll P. Vance. Legumes: Importance and Constraints to Greater Use. \emph{PLANT PHYSIOLOGY}, 2003. DOI:10.1104/pp.017004
\bibitem{ref4} Roger A. Sheldon, John M. Woodley. Role of Biocatalysis in Sustainable Chemistry. \emph{Chemical Reviews}, 2017. DOI:10.1021/acs.chemrev.7b00203
\bibitem{ref5} Roger A. Sheldon. Green solvents for sustainable organic synthesis: state of the art. \emph{Green Chemistry}, 2005. DOI:10.1039/b418069k
\bibitem{ref6} Amar K. Mohanty. Natural Fibers, Biopolymers, and Biocomposites. \emph{}, 2005. DOI:10.1201/9780203508206
\bibitem{ref7} Chidambaram Gunanathan, David Milstein. Applications of Acceptorless Dehydrogenation and Related Transformations in Chemical Synthesis. \emph{Science}, 2013. DOI:10.1126/science.1229712
\bibitem{ref8} Roger A. Sheldon. Fundamentals of green chemistry: efficiency in reaction design. \emph{Chemical Society Reviews}, 2011. DOI:10.1039/c1cs15219j
\end{thebibliography}
\end{document}
